\subsection{Security and Efficiency Trade-Offs}\label{sec:security_of_cryptographic_algorithms:trade_offs}

TTT
%Moreover, it may not always be easy to \textbf{balance security and efficiency} in the implementation and use of cryptographic algorithms due to their interplay: strong security measures must be evaluated against their computational and operational overhead, while optimized implementations must be analyzed for potential weaknesses. This interplay underscores the importance of systematically ensuring and evaluating both the security and efficiency of cryptographic implementations to build resilient and performant systems.


As is often the case in many applications and domains (e.g., \cite{centenaro_safety_2021,berlato_multi_2024}), implementing cryptographic algorithms requires considering several trade-offs between security and efficiency. Typically, strong security measures impose efficiency costs, while optimizing code for efficiency introduces security vulnerabilities (although not necessarily, e.g., see bitslicing in \Cref{sec:security_of_cryptographic_algorithms:best_practices:bitslicing}). Below, we discuss the most relevant of these trade-offs by focusing on cryptographic algorithm implementation --- and not on design or configuration choices such as the choice of cryptographic key lengths.


\subsubsection{Constant-Time Cryptography and Optimizations}\label{sec:security_of_cryptographic_algorithms:trade_offs:constant_time_and_optimization}

One of the clearest examples of security and efficiency trade-offs is constant-time cryptography (security --- see \Cref{sec:security_of_cryptographic_algorithms:best_practices:constant_time}) versus optimizations (efficiency). In fact, making the execution time of instructions constant regardless of their input prevents developers from applying efficiency optimizations like conditional branches for early exits.


\subsubsection{Noise Injection and Performance}\label{sec:security_of_cryptographic_algorithms:trade_offs:noise_injection_and_performance}

Noise injection lowers the \gls{snr} and improves robustness against side-channel attacks such as power analysis (security --- see Note 2.2.3 in \Cref{sec:security_of_cryptographic_algorithms:attacks:sca}) while, however, worsening performance (efficiency). In fact, employing methods like masking, blinding techniques, and dummy procedures (see \Cref{sec:security_of_cryptographic_algorithms:best_practices}) requires executing additional operations (e.g., generating and processing random values), increasing latency and computational overhead.


\subsubsection{Pre-Computation and Timing Attacks}\label{sec:security_of_cryptographic_algorithms:trade_offs:pre_computation_and_timing_attacks}

Pre-computation (efficiency --- see \Cref{sec:security_of_cryptographic_algorithms:performance:best_practices}) may introduce timing attacks (security --- see \Cref{sec:security_of_cryptographic_algorithms:attacks:sca}). In fact, relying on pre-computed values such as lookup tables (e.g., S-boxes in AES), although preventing redundant computations and reducing latency, often introduces memory access patterns that may be exploitable by an attacker to infer (portion of) secret data.


\subsubsection{Memory Management and Performance}\label{sec:security_of_cryptographic_algorithms:trade_offs:memory_management_and_performance}

Proper memory management (security) prevents secret data leakage and memory-related vulnerabilities such as buffer overflows while impacting performance (efficiency). In fact, techniques like memory zeroization and memory inspection (see \Cref{sec:security_of_cryptographic_algorithms:best_practices:zeroization,sec:security_of_cryptographic_algorithms:best_practices:safety}) require additional checks yielding runtime overhead.


\subsubsection{Compilers and Interpreters and Optimizations}\label{sec:security_of_cryptographic_algorithms:trade_offs:compilers_and_interpreters_and_optimizations}

As already discussed in \Cref{sec:security_of_cryptographic_algorithms:languages_and_compilers:compilers,sec:security_of_cryptographic_algorithms:performance:languages_and_compilers:compilers}, compilers and interpreters can greatly optimize the efficiency of cryptographic algorithm implementations (efficiency) while, however, introducing vulnerabilities and even nullifying security measures \cite{schneider_breaking_2024} (security).


\subsubsection{Error Handling and Performance}\label{sec:security_of_cryptographic_algorithms:trade_offs:error_handling_and_performance}

Handling runtime errors without leaking sensitive information (security) may sometimes add computational overhead (efficiency). In fact, preventing unintended behavior from uncaught or mishandled exceptions and signals (see \Cref{sec:security_of_cryptographic_algorithms:languages_and_compilers:languages}) requires developing additional (and possibly complex) logic within cryptographic algorithm implementation.










\sbnote{MAY be RELEVANT FOR SIDE-CHANNEL \url{https://eprint.iacr.org/2025/1130.pdf}}